\section{The weighted range at its right end}\label{sec:boundary}

Theorem C is stated for $p<p_1(w)$, an open condition in $p$ for a fixed exponent $\sigma$. This
section shows that the end point $p=p_1(w)$ is included as well, because the reverse H\"older
exponent of an $A_1$ weight can always be raised. The argument is the classical self-improvement
of reverse H\"older inequalities \cite{CF74}, written out with explicit constants; it uses
nothing about Bochner--Riesz means beyond Theorem C, and adds nothing at the corner $P_2$ of
Theorem B.

In this section cubes have sides parallel to the axes. Let $\omega_n=|B(0,1)|$ and
\[
\kappa_n=\frac{\omega_n n^{n/2}}{2^n},\qquad \mu_n=\frac{2^n}{\omega_n},
\]
the ratio of the volume of a ball to that of the cube inscribed in it, and of a cube to that of
the ball inscribed in it.

\begin{lemma}\label{lem:gehring}
Let $v$ be a weight and $K\ge1$ with $\fint_Qv\le K\operatorname*{ess\,inf}_Qv$ for every cube $Q$. Then
for $0<\eta\le(2^{n+1}K)^{-1}$ and every cube $Q$,
\[
\fint_Qv^{1+\eta}\le2\Bigl(\fint_Qv\Bigr)^{1+\eta}.
\]
\end{lemma}

\begin{proof}
Fix $Q$ and put $b=\fint_Qv>0$. For $x\in Q$ let $M_Qv(x)$ be the supremum of $\fint_Pv$ over the
dyadic subcubes $P$ of $Q$ containing $x$. Then $M_Qv\ge b$ on $Q$, and $v\le M_Qv$ almost
everywhere on $Q$ by the Lebesgue differentiation theorem along dyadic cubes. For each dyadic
subcube $P$ the hypothesis gives $\fint_Pv\le Kv(x)$ for almost every $x\in P$; as there are
countably many such $P$, $M_Qv\le Kv$ almost everywhere on $Q$.

For $\lambda\ge b$ let $\Omega_\lambda=\{x\in Q:M_Qv(x)>\lambda\}$. It is the disjoint union of the
maximal dyadic subcubes $Q_j$ of $Q$ with $\fint_{Q_j}v>\lambda$. None of them is $Q$, since
$\fint_Qv=b\le\lambda$, and by maximality the parent $\widehat Q_j$ of $Q_j$ has
$\fint_{\widehat Q_j}v\le\lambda$. Hence
\begin{equation}\label{eq:gehringlevel}
v(\Omega_\lambda)=\sum_j\int_{Q_j}v\le\sum_j\int_{\widehat Q_j}v\le2^n\lambda\sum_j|Q_j|
=2^n\lambda|\Omega_\lambda|.
\end{equation}

Let $L\ge b$, $F=\min\{M_Qv,L\}$ and $I=\int_QF^\eta v\le L^\eta v(Q)<\infty$. On $Q$ we have
$F\ge b$, and for $b\le\lambda<L$ the set $\{F>\lambda\}$ is $\Omega_\lambda$. The layer-cake
formula and \eqref{eq:gehringlevel} give
\begin{align*}
I&=\int_0^b\eta\lambda^{\eta-1}v(Q)\,d\lambda+\int_b^L\eta\lambda^{\eta-1}v(\Omega_\lambda)\,d\lambda
\le b^\eta v(Q)+2^n\eta\int_0^\infty\lambda^\eta\bigl|\{x\in Q:F(x)>\lambda\}\bigr|\,d\lambda\\
&=b^\eta v(Q)+\frac{2^n\eta}{1+\eta}\int_QF^{1+\eta}
\le b^\eta v(Q)+2^nK\eta\, I,
\end{align*}
since $F^{1+\eta}=F^\eta F\le KF^\eta v$. As $2^nK\eta\le\frac12$ and $I$ is finite,
$I\le2b^\eta v(Q)$. Letting $L\to\infty$, monotone convergence and $v\le M_Qv$ give
$\int_Qv^{1+\eta}\le\int_Q(M_Qv)^\eta v\le2b^\eta v(Q)=2b^{1+\eta}|Q|$.
\end{proof}

\begin{proposition}\label{prop:selfimprove}
Let $\sigma>1$, $w\in A_1\cap RH_\sigma$, $A=[w]_{A_1}$, $B=[w]_{RH_\sigma}$, and
\[
K=\kappa_n(AB)^\sigma,\qquad\eta=\frac1{2^{n+1}K},\qquad\Sigma=\sigma(1+\eta).
\]
Then $w\in RH_\Sigma$ and $[w]_{RH_\Sigma}\le(2\mu_nK^{1+\eta})^{1/\Sigma}B$.
\end{proposition}

\begin{proof}
Put $v=w^\sigma$. For every ball $B'$, as shown after Theorem \ref{thm:fn},
$\fint_{B'}v\le(AB)^\sigma\operatorname*{ess\,inf}_{B'}v$. If $Q$ is a cube and $B'$ the ball circumscribed
about it, then $|B'|=\kappa_n|Q|$ and
\[
\fint_Qv\le\kappa_n\fint_{B'}v\le\kappa_n(AB)^\sigma\operatorname*{ess\,inf}_{B'}v
\le K\operatorname*{ess\,inf}_Qv .
\]
Here $K\ge1$, since $A\ge1$, $B\ge1$ and $\kappa_n\ge1$, so Lemma \ref{lem:gehring} applies to $v$.
Now let $B'$ be a ball and $Q$ the cube circumscribed about it, so $|Q|=\mu_n|B'|$. By the lemma
and the bound just proved,
\[
\fint_{B'}v^{1+\eta}\le\mu_n\fint_Qv^{1+\eta}\le2\mu_n\Bigl(\fint_Qv\Bigr)^{1+\eta}
\le2\mu_n\bigl(K\operatorname*{ess\,inf}_Qv\bigr)^{1+\eta}\le2\mu_nK^{1+\eta}\Bigl(\fint_{B'}v\Bigr)^{1+\eta},
\]
since $\operatorname*{ess\,inf}_Qv\le\operatorname*{ess\,inf}_{B'}v\le\fint_{B'}v$. With $v=w^\sigma$ and
$\fint_{B'}w^\sigma\le B^\sigma(\fint_{B'}w)^\sigma$ this reads
$\fint_{B'}w^\Sigma\le2\mu_nK^{1+\eta}B^\Sigma(\fint_{B'}w)^\Sigma$.
\end{proof}

\begin{theorem}\label{thm:boundary}
Let $n\ge2$, $a>0$, $\sigma>1$ and $w\in A_1\cap RH_\sigma$, and let $\Sigma>\sigma$ be the exponent of
Proposition \ref{prop:selfimprove}. Then the conclusions of Theorem C hold for
\[
1<p<1+\frac{n-1}{n+1}\Bigl(1-\frac1\Sigma\Bigr),
\]
and in particular at $p=p_1(w)=1+\frac{n-1}{n+1}(1-\frac1\sigma)$. Equivalently, with
$\sigma_w=\sup\{s>1:w\in RH_s\}\in(1,\infty]$, they hold for
$1<p<1+\frac{n-1}{n+1}(1-\frac1{\sigma_w})$, and $w\notin RH_{\sigma_w}$ when $\sigma_w<\infty$.
\end{theorem}

\begin{proof}
By Proposition \ref{prop:selfimprove}, $w\in A_1\cap RH_\Sigma$, and Theorem C with $\Sigma$ in place
of $\sigma$ gives the first assertion; since $s\mapsto1-\frac1s$ is increasing and $\Sigma>\sigma$,
the range contains $p_1(w)$. The set of $s>1$ with $w\in RH_s$ is an interval by H\"older's
inequality, nonempty by \cite{CF74}, and open by Proposition \ref{prop:selfimprove}; so its
supremum $\sigma_w$ is not attained, and every $p<1+\frac{n-1}{n+1}(1-\frac1{\sigma_w})$ is below
$p_1$ for some admissible exponent.
\end{proof}

\begin{corollary}\label{cor:A1only}
Let $n\ge2$, $a>0$ and $w\in A_1$. The conclusions of Theorem C hold for
\[
1<p<1+\frac{n-1}{(n+1)\bigl(2^{n+1}\kappa_n[w]_{A_1}+1\bigr)} .
\]
\end{corollary}

\begin{proof}
The second half of the proof of Proposition \ref{prop:selfimprove}, applied to $v=w$ with
$K=\kappa_n[w]_{A_1}$ in place of $\kappa_n(AB)^\sigma$, shows that $w\in RH_{1+\eta_0}$ with
$\eta_0=(2^{n+1}\kappa_n[w]_{A_1})^{-1}$. Theorem C with
$\sigma=1+\eta_0$ gives the range, since $1-\frac1{1+\eta_0}=\frac{\eta_0}{1+\eta_0}
=\frac1{2^{n+1}\kappa_n[w]_{A_1}+1}$.
\end{proof}

\begin{example}\label{ex:powerweight}
For $0\le\beta<n$ the weight $|x|^{-\beta}$ lies in $A_1$ \cite[\S7.1]{Gra14}. Let $0\le\gamma<n$ and
$w(x)=|x|^{-\gamma}$. If $\gamma s<n$, then $w^s\in A_1$, so
$\fint_Bw^s\le C\operatorname*{ess\,inf}_Bw^s\le C(\fint_Bw)^s$ for every ball and $w\in RH_s$; if
$\gamma s\ge n$, then $w^s$ is not integrable near $0$ and $w\notin RH_s$. Thus $\sigma_w=n/\gamma$, and
Theorem \ref{thm:boundary} gives the weak-type bound on $L^p(|x|^{-\gamma})$ for
\[
1<p<1+\frac{n-1}{n+1}\Bigl(1-\frac\gamma n\Bigr).
\]
For $\gamma=0$ this is Theorem A.
\end{example}
